% =========================================================================
% Capítulo 3: Cinética 2 - Dinámica Avanzada de Sistemas de Partículas
% =========================================================================
% Archivo: C3_00_introduccion.tex
% =========================================================================

\chapter{Cinética 2: Sistemas Avanzados y Cuerpos Deformables}

\section*{Introducción}

En este capítulo ampliamos el análisis dinámico a sistemas más complejos: sistemas con masa variable, sistemas continuos, y los efectos de fuerzas internas en la dinámica de sistemas extendidos.

\subsection*{Temas Principales}

\begin{enumerate}
    \item \textbf{Sistemas de Masa Variable}: Casos como cohetes que expelen masa.
    \item \textbf{Sistemas Continuos}: Tratamiento de distribuciones continuas de masa.
    \item \textbf{Vibraciones}: Movimientos periódicos en sistemas acoplados.
    \item \textbf{Interacciones Múltiples}: Sistemas con muchos cuerpos interactuando.
\end{enumerate}

\section{Ecuaciones de Movimiento Avanzadas}

\subsection*{Sistemas de Masa Variable}

Para un sistema cuya masa cambia con el tiempo:

$$\frac{d(m\vec{v})}{dt} = \vec{F}_{\text{ext}} + \vec{v}_{\text{rel}} \frac{dm}{dt}$$

donde $\vec{v}_{\text{rel}}$ es la velocidad relativa de la masa que se añade o se retira.

\subsection*{Energía en Sistemas Complejos}

\begin{description}
    \item[Energía Cinética Absoluta:] Suma de energías cinéticas de todas las partículas.
    \item[Energía Potencial:] Asociada con fuerzas conservativas (gravitatoria, elástica, etc.).
    \item[Disipación:] Trabajo realizado contra fuerzas disipativas (fricción, viscosidad).
    \item[Trabajo Virtual:] Concepto clave en sistemas con restricciones.
\end{description}

\subsection*{Teorema de Trabajo-Energía}

$$\Delta T = W_{\text{neto}} = W_{\text{ext}} + W_{\text{int}}$$

Para sistemas con restricciones, el trabajo de fuerzas internas es especialmente importante.

\section{Análisis de Vibraciones}

\subsection*{Oscilador Armónico Simple}

$$\ddot{x} + \omega_0^2 x = 0$$

donde $\omega_0 = \sqrt{k/m}$ es la frecuencia natural.

\subsection*{Ecuación de Lagrange de Segunda Especie}

$$\frac{d}{dt}\left(\frac{\partial L}{\partial \dot{q}_j}\right) - \frac{\partial L}{\partial q_j} = Q_j$$

donde $L = T - V$ es el Lagrangiano, $q_j$ son coordenadas generalizadas y $Q_j$ son fuerzas generalizadas.

\section*{Objetivos del Capítulo}

Al completar este capítulo, deberás:
\begin{itemize}
    \item Analizar sistemas dinámicos complejos y acoplados.
    \item Manejar sistemas de masa variable.
    \item Entender la energía en sistemas disipativos.
    \item Aplicar la ecuación de Lagrange.
    \item Resolver problemas de oscilaciones y vibraciones.
\end{itemize}

\newpage
